% GENERATED FILE. Edit canonical lessons, then run npm run generate:books.
% canonical-source-hash: fnv1a64:9f2352c00c6a548b
% canonical-lessons: TA-C142-ottakam, TA-C142-vattu, TA-C142-kannam, TA-C142-manikkattu, TA-C142-mulankai

\chapter{Camel, Duck, Cheek, Wrist, Elbow}
\label{ch:ta-camel-duck-cheek-wrist-elbow}
\hlchaptermodality{142}

\begin{chapteropening}
\textbf{By the end of this chapter:} \emph{I can say a camel, a duck, a cheek, a wrist and an elbow.}

Complete the last lesson of chapter 142: I can say a camel, a duck, a cheek, a wrist and an elbow.
\end{chapteropening}

\section[\texorpdfstring{\ta{ஒட்டகம்} (ottakam)}{ஒட்டகம் (ottakam)}]{\ta{ஒட்டகம்} (ottakam) — a camel}

\label{lesson:TA-C142-ottakam}



\begin{quote}
Before the new one: say the Tamil for a pig, then the Tamil for a donkey.
\end{quote}

\subsection*{You'll want to know: \ta{ஒட்டகம்}}
\textbf{\ta{ஒட்டகம்}} — \emph{ottakam} — \textquotedblleft{}a camel\textquotedblright{}.

\textbf{In the family, and next door}

\noindent
\begin{tabularx}{\linewidth}{@{}>{\raggedright\arraybackslash}X>{\raggedright\arraybackslash}X>{\raggedright\arraybackslash}X@{}}
\toprule
\textbf{} & \textbf{word} & \textbf{same root} \\
\midrule
Kannada & \ka{ಒಂಟೆ} (\emph{oṇṭe}) & yes \\
Telugu & \te{ఒంటె} (\emph{oṇṭe}) & yes \\
Malayalam & \ml{ഒട്ടകം} (\emph{oṭṭakaṁ}) & yes \\
Hindi (neighbour) & \dv{ऊँट} (\emph{ū̃ṭ}) &  \\
English & a camel &  \\
\bottomrule
\end{tabularx}

\begin{grammarlens}[title={What you've built}]
One more word for this chapter.
\end{grammarlens}

\subsection*{Your turn}
\begin{itemize}
\raggedright
  \item \emph{Say it:} \emph{ottakam}
  \item \emph{Say it:} \emph{ottakam}, once more
  \item \emph{Say it:} say \emph{ottakam}
  \item \emph{Recall:} say the Tamil for a pig, then the Tamil for a donkey, then say \emph{ottakam} again
\end{itemize}

\begin{tcolorbox}[breakable,colback=teal!4,colframe=teal!35!black,title={Before you move on}]
What does \ta{ஒட்டகம்} mean? (\textquotedblleft{}A camel\textquotedblright{}.) Say it once more.
\end{tcolorbox}
\section[\texorpdfstring{\ta{வாத்து} (vāttu)}{வாத்து (vāttu)}]{\ta{வாத்து} (vāttu) — a duck}

\label{lesson:TA-C142-vattu}



\begin{quote}
Before the new one: say the Tamil for a donkey, then the Tamil for a camel.
\end{quote}

\subsection*{You'll want to know: \ta{வாத்து}}
\textbf{\ta{வாத்து}} — \emph{vāttu} — \textquotedblleft{}a duck\textquotedblright{}.

\textbf{In the family, and next door}

\noindent
\begin{tabularx}{\linewidth}{@{}>{\raggedright\arraybackslash}X>{\raggedright\arraybackslash}X@{}}
\toprule
\textbf{} & \textbf{word} \\
\midrule
Kannada & \ka{ಬಾತುಕೋಳಿ} (\emph{bātukōḷi}) \\
Telugu & \te{బాతు} (\emph{bātu}) \\
Malayalam & \ml{താറാവ്} (\emph{tāṟāvŭ}) \\
English & a duck \\
\bottomrule
\end{tabularx}

\begin{grammarlens}[title={What you've built}]
One more word for this chapter.
\end{grammarlens}

\subsection*{Your turn}
\begin{itemize}
\raggedright
  \item \emph{Say it:} \emph{vāttu}
  \item \emph{Say it:} \emph{vāttu}, once more
  \item \emph{Say it:} say \emph{vāttu}
  \item \emph{Recall:} say the Tamil for a donkey, then the Tamil for a camel, then say \emph{vāttu} again
\end{itemize}

\begin{tcolorbox}[breakable,colback=teal!4,colframe=teal!35!black,title={Before you move on}]
What does \ta{வாத்து} mean? (\textquotedblleft{}A duck\textquotedblright{}.) Say it once more.
\end{tcolorbox}
\section[\texorpdfstring{\ta{கன்னம்} (kaṉṉam)}{கன்னம் (kaṉṉam)}]{\ta{கன்னம்} (kaṉṉam) — a cheek}

\label{lesson:TA-C142-kannam}



\begin{quote}
Before the new one: say the Tamil for a camel, then the Tamil for a duck.
\end{quote}

\subsection*{You'll want to know: \ta{கன்னம்}}
\textbf{\ta{கன்னம்}} — \emph{kaṉṉam} — \textquotedblleft{}a cheek\textquotedblright{}.

\textbf{In the family, and next door}

\noindent
\begin{tabularx}{\linewidth}{@{}>{\raggedright\arraybackslash}X>{\raggedright\arraybackslash}X>{\raggedright\arraybackslash}X@{}}
\toprule
\textbf{} & \textbf{word} & \textbf{same root} \\
\midrule
Kannada & \ka{ಕೆನ್ನೆ} (\emph{kenne}) & yes \\
Telugu & \te{బుగ్గ} (\emph{bugga}) &  \\
Malayalam & \ml{കവിൾ} (\emph{kavil}) &  \\
Hindi (neighbour) & \dv{गाल} (\emph{gāl}) &  \\
English & a cheek &  \\
\bottomrule
\end{tabularx}

\begin{grammarlens}[title={What you've built}]
One more word for this chapter.
\end{grammarlens}

\subsection*{Your turn}
\begin{itemize}
\raggedright
  \item \emph{Say it:} \emph{kaṉṉam}
  \item \emph{Say it:} \emph{kaṉṉam}, once more
  \item \emph{Say it:} say \emph{kaṉṉam}
  \item \emph{Recall:} say the Tamil for a camel, then the Tamil for a duck, then say \emph{kaṉṉam} again
\end{itemize}

\begin{tcolorbox}[breakable,colback=teal!4,colframe=teal!35!black,title={Before you move on}]
What does \ta{கன்னம்} mean? (\textquotedblleft{}A cheek\textquotedblright{}.) Say it once more.
\end{tcolorbox}
\section[\texorpdfstring{\ta{மணிக்கட்டு} (maṇikkaṭṭu)}{மணிக்கட்டு (maṇikkaṭṭu)}]{\ta{மணிக்கட்டு} (maṇikkaṭṭu) — a wrist}

\label{lesson:TA-C142-manikkattu}



\begin{quote}
Before the new one: say the Tamil for a duck, then the Tamil for a cheek.
\end{quote}

\subsection*{You'll want to know: \ta{மணிக்கட்டு}}
\textbf{\ta{மணிக்கட்டு}} — \emph{maṇikkaṭṭu} — \textquotedblleft{}a wrist\textquotedblright{}.

\textbf{In the family, and next door}

\noindent
\begin{tabularx}{\linewidth}{@{}>{\raggedright\arraybackslash}X>{\raggedright\arraybackslash}X>{\raggedright\arraybackslash}X@{}}
\toprule
\textbf{} & \textbf{word} & \textbf{same root} \\
\midrule
Kannada & \ka{ಮಣಿಕಟ್ಟು} (\emph{maṇikaṭṭu}) & yes \\
Telugu & \te{మణికట్టు} (\emph{maṇikaṭṭu}) & yes \\
Hindi (neighbour) & \dv{कलाई} (\emph{kalāī}) &  \\
English & a wrist &  \\
\bottomrule
\end{tabularx}

\begin{grammarlens}[title={What you've built}]
One more word for this chapter.
\end{grammarlens}

\subsection*{Your turn}
\begin{itemize}
\raggedright
  \item \emph{Say it:} \emph{maṇikkaṭṭu}
  \item \emph{Say it:} \emph{maṇikkaṭṭu}, once more
  \item \emph{Say it:} say \emph{maṇikkaṭṭu}
  \item \emph{Recall:} say the Tamil for a duck, then the Tamil for a cheek, then say \emph{maṇikkaṭṭu} again
\end{itemize}

\begin{tcolorbox}[breakable,colback=teal!4,colframe=teal!35!black,title={Before you move on}]
What does \ta{மணிக்கட்டு} mean? (\textquotedblleft{}A wrist\textquotedblright{}.) Say it once more.
\end{tcolorbox}
\section[\texorpdfstring{\ta{முழங்கை} (muḻaṅkai)}{முழங்கை (muḻaṅkai)}]{\ta{முழங்கை} (muḻaṅkai) — an elbow}

\label{lesson:TA-C142-mulankai}



\begin{quote}
Before the new one: say the Tamil for a cheek, then the Tamil for a wrist.
\end{quote}

\subsection*{You'll want to know: \ta{முழங்கை}}
\textbf{\ta{முழங்கை}} — \emph{muḻaṅkai} — \textquotedblleft{}an elbow\textquotedblright{}.

\textbf{In the family, and next door}

\noindent
\begin{tabularx}{\linewidth}{@{}>{\raggedright\arraybackslash}X>{\raggedright\arraybackslash}X>{\raggedright\arraybackslash}X@{}}
\toprule
\textbf{} & \textbf{word} & \textbf{same root} \\
\midrule
Kannada & \ka{ಮೊಣಕೈ} (\emph{moṇakai}) & yes \\
Telugu & \te{మోచేయి} (\emph{mōcēyi}) &  \\
Malayalam & \ml{കൈമുട്ട്} (\emph{kaimuṭṭŭ}) &  \\
Hindi (neighbour) & \dv{कोहनी} (\emph{kohnī}) &  \\
English & an elbow &  \\
\bottomrule
\end{tabularx}

\begin{grammarlens}[title={What you've built}]
One more word for this chapter.
\end{grammarlens}

\subsection*{Your turn}
\begin{itemize}
\raggedright
  \item \emph{Say it:} \emph{muḻaṅkai}
  \item \emph{Say it:} \emph{muḻaṅkai}, once more
  \item \emph{Say it:} say \emph{muḻaṅkai}
  \item \emph{Recall:} say the Tamil for a cheek, then the Tamil for a wrist, then say \emph{muḻaṅkai} again
\end{itemize}

\begin{tcolorbox}[breakable,colback=teal!4,colframe=teal!35!black,title={Before you move on}]
What does \ta{முழங்கை} mean? (\textquotedblleft{}An elbow\textquotedblright{}.) Say it once more.
\end{tcolorbox}
